\begin{pastexam}[label={px:ncku-114-prog-15}]{成大資訊聯招 114 學年度 第 15 題}
（填空）

(10\%) We offer $n$ courses where each course $i$ has a start time $s_i$, a finish time $f_i$, and a weight $w_i$. Two courses are considered \textbf{compatible} if their time intervals do not overlap. For simplicity, we can assume that the courses are sorted in non-decreasing order of their finish times, denoted as $f_1\le f_2\le\cdots\le f_n$. Additionally, let $p(j)$ for a course $j$ represent the largest index $i<j$ such that courses $i$ and $j$ are disjoint. We define $p(j)=0$ if no course $i<j$ is disjoint from $j$. Our goal is to find a set $S\subseteq\{1,\ldots,n\}$ of mutually compatible courses such that the total weight of the courses in $S$ is maximized. We have designed a recursive algorithm for this purpose, but we have omitted an essential part. Please assist us in completing the missing portion. \textbf{Hint:} This algorithm returns the value of the optimal solution to the problem consisting of courses $\{1,\ldots,i\}$.
\begin{center}
\begin{tabular}{|p{0.85\textwidth}|}
\hline
F($i$) \\
\quad If $i=0$ return 0 \\
\quad return\underline{\qquad (10\%)\qquad} \\
\hline
\end{tabular}
\end{center}
\end{pastexam}

\begin{pastsoln}
答案：$\max\bigl(w_i + F(p(i)),\ F(i - 1)\bigr)$。

這是加權區間排程（weighted interval scheduling），資結書第~19.1~節的活動選擇加上權重以後，貪心不再成立，要用動態規劃
（資結書第~18.1.5~節的步驟）。 $F(i)$ 是只看課程 $1, \ldots, i$ 時的最佳總權重。 考慮最佳解裡有沒有課程 $i$：
\begin{itemize}
\item \textbf{選了 $i$}：和 $i$ 重疊的課都不能選。 課程依結束時間排好，$p(i)$ 是和 $i$ 不重疊的最大編號：$p(i)$ 之後、$i$ 之前的
課結束時間都晚於 $i$ 的開始時間（否則 $p(i)$ 會更大），全都和 $i$ 重疊；而 $1, \ldots, p(i)$ 的結束時間都不晚於 $f_{p(i)} \le s_i$，
全部和 $i$ 相容。 所以剩下的就是課程 $1, \ldots, p(i)$ 的子問題，最佳值 $w_i + F(p(i))$。
\item \textbf{沒選 $i$}：就是課程 $1, \ldots, i - 1$ 的子問題，最佳值 $F(i - 1)$。
\end{itemize}
兩者取大：
\[
F(i) = \begin{cases} 0, & i = 0, \\ \max\bigl(w_i + F(p(i)),\ F(i - 1)\bigr), & i \ge 1. \end{cases}
\]
空格填 \texttt{max(w[i] + F(p(i)), F(i - 1))}。 照題目寫成純遞迴會重複計算同樣的子問題，最壞是指數時間；把 $F(0..n)$ 存在
表格裡由小到大填（或加備忘錄），每格 $O(1)$，加上算 $p(\cdot)$ 的 $O(n \log n)$（二分搜尋），總共 $O(n \log n)$。

小例子核對：三門課 $[0, 3)$ 權重 $2$、$[2, 5)$ 權重 $4$、$[4, 7)$ 權重 $4$，$p(1) = 0$、$p(2) = 0$、$p(3) = 1$。
$F(1) = 2$，$F(2) = \max(4 + 0, 2) = 4$，$F(3) = \max(4 + F(1), F(2)) = 6$，選課程 $1$、$3$，正確。
\end{pastsoln}
